In this section, we present a representative case comparing three model variants: \textbf{GraphWalker-SFT-GraphSynth}, \textbf{GraphWalker-SFT-GraphRoll}, and \textbf{GraphWalker-SFT-RL}. We analyze their reasoning trajectories on one question from CWQ: \textit{``Where is the home stadium of the team who won the 1946 World Series championship?''}
As shown in Appendix~\ref{app:four_cases}, \textbf{GraphWalker-SFT-GraphSynth} successfully retrieves the correct triple \texttt{[St. Louis Cardinals, arena\_stadium, Busch Stadium]} at Turn 3, and even further verifies it via an additional query at Turn 4, explicitly confirming \textit{``Busch Stadium is the home venue for the St. Louis Cardinals''} in its reasoning, yet it outputs \texttt{[``St. Louis Cardinals'']} instead of extracting the stadium name---this reveals that Stage 1 SFT (on GraphSynth) teaches KG navigation but not answer extraction.
\textbf{GraphWalker-SFT-GraphRoll} identifies the champion St. Louis Cardinals at Turn 2 and observes \texttt{Sportsman's Park} in the event locations, but directly outputs this historical venue \texttt{[``Sportsman's Park'']} without further exploring the team's current stadium relation---this shows that Stage 2 SFT (on GraphRoll) alone, trained on short expert trajectories, lacks the multi-hop exploration capacity needed for deeper reasoning paths.
\textbf{GraphWalker-SFT-RL} correctly identifies the champion at Turn 2, then directly queries the St. Louis Cardinals' stadium relation and outputs \texttt{[``Busch Stadium'']} as the final answer (EM=1.0), demonstrating that RL consolidates the exploration capacity of Stage 1 and the reflection capability of Stage 2 into a more efficient and accurate reasoning strategy.


\subsection{Complete Reasoning Trajectories}
\label{app:four_cases}

We present the complete reasoning trajectories of all four models. Each trajectory is displayed in a table where different message types are color-coded: \colorbox{blue!10}{thinking}, \colorbox{green!10}{KG queries}, \colorbox{yellow!10}{information feedback}, \colorbox{red!10}{errors}, and final answers.

\subsubsection{GraphWalker-SFT-GraphSynth}

{\scriptsize
\setlength{\arrayrulewidth}{0.3pt}
\setlength{\tabcolsep}{8pt}
\renewcommand{\arraystretch}{1.3}
\begin{longtable}{>{\raggedright\arraybackslash}p{0.97\textwidth}}
\toprule
\rowcolor{gray!20}
\textbf{Question:} Where is the home stadium of the team who won the 1946 World Series championship? \\
\rowcolor{gray!20}
\textbf{Initial Entities:} 1946 World Series \\
\midrule
\rowcolor{blue!10}
\texttt{<think>} Since I need to determine the home stadium of the team that won the 1946 World Series, I will start by querying the relations for the entity "1946 World Series". This will help me identify relevant connections to the winning team and its stadium. Gathering these relations is essential for finding the necessary information to answer the question. \texttt{</think>} \\
\rowcolor{green!10}
\texttt{<kg-query>} get\_relations("1946 World Series") \texttt{</kg-query>} \\
\rowcolor{yellow!10}
\texttt{<information>} sports.sports\_championship\_event.champion, sports.sports\_championship\_event.result, sports.sports\_team.championships, time.event.locations, ... \texttt{</information>} \\
\rowcolor{blue!10}
\texttt{<think>} The relations obtained for "1946 World Series" include "sports.sports\_championship\_event.champion," which is crucial for identifying the winning team, and "time.event.locations," which may provide information about the home stadium. By querying the triples with these relations, I can gather detailed information about the champion team and its associated stadium, which is necessary to answer the question accurately. \texttt{</think>} \\
\rowcolor{green!10}
\texttt{<kg-query>} get\_triples("1946 World Series", ["sports.sports\_championship\_event.champion", "sports.sports\_championship\_event.result", "sports.sports\_team.championships", "time.event.locations", ...]) \texttt{</kg-query>} \\
\rowcolor{yellow!10}
\texttt{<information>} [1946 World Series, sports.sports\_championship\_event.champion, St. Louis Cardinals], [1946 World Series, sports.sports\_championship\_event.result, 4-3], [1946 World Series, sports.sports\_championship\_event.championship, World Series], [St. Louis Cardinals, sports.sports\_team.championships, 1946 World Series] \texttt{</information>} \\
\rowcolor{blue!10}
\texttt{<think>} I learned that the champion of the 1946 World Series is the St. Louis Cardinals. To find the home stadium of this team, I need to query the relations for "St. Louis Cardinals," as this will provide insights into their stadium and other relevant details. This step is essential to connect the team with its home venue, which is necessary to answer the question accurately. \texttt{</think>} \\
\rowcolor{green!10}
\texttt{<kg-query>} get\_relations("St. Louis Cardinals") \texttt{</kg-query>} \\
\rowcolor{yellow!10}
\texttt{<information>} sports.sports\_team.venue, sports.sports\_team.arena\_stadium, sports.sports\_team.location, ... \texttt{</information>} \\
\rowcolor{blue!10}
\texttt{<think>} The relations obtained for "St. Louis Cardinals" include "sports.sports\_team.venue" and "sports.sports\_team.arena\_stadium," which are directly relevant for identifying the team's home stadium. By querying the triples using these relations, I can gather the specific information needed to answer the question about where the St. Louis Cardinals play their home games. \texttt{</think>} \\
\rowcolor{green!10}
\texttt{<kg-query>} get\_triples("St. Louis Cardinals", ["sports.sports\_team.venue", "sports.sports\_team.arena\_stadium", "sports.sports\_team.location", ...]) \texttt{</kg-query>} \\
\rowcolor{yellow!10}
\texttt{<information>} [St. Louis Cardinals, sports.sports\_team.arena\_stadium, Busch Stadium], [St. Louis Cardinals, sports.sports\_team.arena\_stadium, Roger Dean Stadium], [St. Louis Cardinals, sports.sports\_team.location, St. Louis], ... \texttt{</information>} \\
\rowcolor{blue!10}
\texttt{<think>} The last observation revealed that the St. Louis Cardinals have multiple venues listed, including "Busch Stadium." However, to confirm the current home stadium, I need to further investigate the relations associated with "Busch Stadium." This will help clarify if it is indeed the primary venue used by the team for their home games, which is crucial for answering the question accurately. \texttt{</think>} \\
\rowcolor{green!10}
\texttt{<kg-query>} get\_relations("Busch Stadium") \texttt{</kg-query>} \\
\rowcolor{yellow!10}
\texttt{<information>} sports.sports\_facility.home\_venue\_for, sports.sports\_facility.teams, sports.sports\_team.arena\_stadium, ... \texttt{</information>} \\
\rowcolor{blue!10}
\texttt{<think>} The relations obtained for "Busch Stadium" include "sports.sports\_facility.home\_venue\_for" and "sports.sports\_facility.teams," which are critical for confirming whether it serves as the home stadium for the St. Louis Cardinals. By querying the triples with these relations, I can verify the connection between the stadium and the team, ensuring that I accurately identify the home stadium for the team that won the 1946 World Series. \texttt{</think>} \\
\rowcolor{green!10}
\texttt{<kg-query>} get\_triples("Busch Stadium", ["sports.sports\_facility.home\_venue\_for", "sports.sports\_facility.teams", "sports.sports\_team.arena\_stadium", ...]) \texttt{</kg-query>} \\
\rowcolor{yellow!10}
\texttt{<information>} [Busch Stadium, sports.sports\_facility.teams, St. Louis Cardinals], [St. Louis Cardinals, sports.sports\_team.arena\_stadium, Busch Stadium], ... \texttt{</information>} \\
\rowcolor{blue!10}
\texttt{<think>} The last observation confirmed that Busch Stadium is indeed the home venue for the St. Louis Cardinals, as indicated by the triple showing "[Busch Stadium, sports.sports\_facility.teams, St. Louis Cardinals]." Since the question asks for the home stadium of the team that won the 1946 World Series, I can confidently conclude that Busch Stadium is the correct answer. Thus, I will provide the answer now. \texttt{</think>} \\
\rowcolor{red!10}
\texttt{<answer>} ["St. Louis Cardinals"] \texttt{</answer>} \textbf{(EM=0.0 - outputs team name instead of stadium)} \\
\bottomrule
\end{longtable}
}


\newpage
\subsubsection{GraphWalker-SFT-GraphRoll}

{\scriptsize
\setlength{\arrayrulewidth}{0.3pt}
\setlength{\tabcolsep}{8pt}
\renewcommand{\arraystretch}{1.3}
\begin{longtable}{>{\raggedright\arraybackslash}p{0.97\textwidth}}
\toprule
\rowcolor{gray!20}
\textbf{Question:} Where is the home stadium of the team who won the 1946 World Series championship? \\
\rowcolor{gray!20}
\textbf{Initial Entities:} 1946 World Series \\
\midrule
\rowcolor{blue!10}
\texttt{<think>} To find the home stadium of the team that won the 1946 World Series championship, I will start by querying the relations of the entity "1946 World Series". This should provide information about the winning team and potentially its home stadium. After retrieving the relations, I will look for relevant details regarding the championship and the winning team. \texttt{</think>} \\
\rowcolor{green!10}
\texttt{<kg-query>} get\_relations("1946 World Series") \texttt{</kg-query>} \\
\rowcolor{yellow!10}
\texttt{<information>} sports.sports\_championship\_event.champion, sports.sports\_championship\_event.result, sports.sports\_team.championships, time.event.locations, ... \texttt{</information>} \\
\rowcolor{blue!10}
\texttt{<think>} The relations retrieved for the "1946 World Series" include information about the champion, result, championships associated with the team, locations, start date, end date, and runner-up. To find the home stadium of the winning team, I will focus on the "champion" relation first, as it should indicate which team won the championship. Once I identify the champion, I can then look for its home stadium. I will proceed to retrieve the triples for the "1946 World Series" using the top relevant relations. \texttt{</think>} \\
\rowcolor{green!10}
\texttt{<kg-query>} get\_triples("1946 World Series", ["sports.sports\_championship\_event.champion", "sports.sports\_championship\_event.result", "sports.sports\_team.championships", "time.event.locations"]) \texttt{</kg-query>} \\
\rowcolor{yellow!10}
\texttt{<information>} [1946 World Series, sports.sports\_championship\_event.champion, St. Louis Cardinals], [1946 World Series, sports.sports\_championship\_event.result, 4-3], [St. Louis Cardinals, sports.sports\_team.championships, 1946 World Series], [1946 World Series, time.event.locations, Fenway Park], [1946 World Series, time.event.locations, Sportsman's Park] \texttt{</information>} \\
\rowcolor{blue!10}
\texttt{<think>} The champion is St. Louis Cardinals. The event locations include Sportsman's Park, which is in St. Louis. Since the question asks for the home stadium of the winning team, and Sportsman's Park appears in the event locations, I will conclude this is the home stadium. \texttt{</think>} \\
\rowcolor{red!10}
\texttt{<answer>} ["Sportsman's Park"] \texttt{</answer>} \textbf{(EM=0.0 - stops after 2 turns without exploring team's stadium relation)} \\
\bottomrule
\end{longtable}
}


\subsubsection{GraphWalker-SFT-RL}

{\scriptsize
\setlength{\arrayrulewidth}{0.3pt}
\setlength{\tabcolsep}{8pt}
\renewcommand{\arraystretch}{1.3}
\begin{longtable}{>{\raggedright\arraybackslash}p{0.97\textwidth}}
\toprule
\rowcolor{gray!20}
\textbf{Question:} Where is the home stadium of the team who won the 1946 World Series championship? \\
\rowcolor{gray!20}
\textbf{Initial Entities:} 1946 World Series \\
\midrule
\rowcolor{blue!10}
\texttt{<think>} To find the home stadium of the team that won the 1946 World Series championship, I will start by querying the relations of the "1946 World Series" entity. This should provide information about the winning team and potentially lead to details about their home stadium. \texttt{</think>} \\
\rowcolor{green!10}
\texttt{<kg-query>} get\_relations("1946 World Series") \texttt{</kg-query>} \\
\rowcolor{yellow!10}
\texttt{<information>} sports.sports\_championship\_event.champion, sports.sports\_championship\_event.result, sports.sports\_championship\_event.runner\_up, ... \texttt{</information>} \\
\rowcolor{blue!10}
\texttt{<think>} The relations returned for the "1946 World Series" include several relevant ones, particularly "champion," which likely indicates the winning team. I will also consider "runner\_up" as it may provide context about the teams involved. To gather more detailed information, I will query the triples for the "1946 World Series" using these relations. \texttt{</think>} \\
\rowcolor{green!10}
\texttt{<kg-query>} get\_triples("1946 World Series", ["sports.sports\_championship\_event.champion", "sports.sports\_championship\_event.runner\_up", "sports.sports\_championship\_event.result"]) \texttt{</kg-query>} \\
\rowcolor{yellow!10}
\texttt{<information>} [1946 World Series, sports.sports\_championship\_event.champion, St. Louis Cardinals], [1946 World Series, sports.sports\_championship\_event.runner\_up, Boston Red Sox], [1946 World Series, sports.sports\_championship\_event.result, 4-3], [1946 World Series, sports.sports\_championship\_event.championship, World Series] \texttt{</information>} \\
\rowcolor{blue!10}
\texttt{<think>} The champion of the 1946 World Series is the St. Louis Cardinals, and the runner-up is the Boston Red Sox. To find the home stadium of the St. Louis Cardinals, I will now query the relations for the "St. Louis Cardinals" entity. This should help me identify the stadium associated with this team. \texttt{</think>} \\
\rowcolor{green!10}
\texttt{<kg-query>} get\_relations("St. Louis Cardinals") \texttt{</kg-query>} \\
\rowcolor{yellow!10}
\texttt{<information>} sports.sports\_team.venue, sports.sports\_team.arena\_stadium, sports.sports\_team.location, ... \texttt{</information>} \\
\rowcolor{blue!10}
\texttt{<think>} The relations for the "St. Louis Cardinals" include "venue" and "arena\_stadium," both of which are likely to provide information about the home stadium of the team. I will query the triples for the "St. Louis Cardinals" using these relations to find the specific stadium. \texttt{</think>} \\
\rowcolor{green!10}
\texttt{<kg-query>} get\_triples("St. Louis Cardinals", ["sports.sports\_team.venue", "sports.sports\_team.arena\_stadium", "sports.sports\_team.location"]) \texttt{</kg-query>} \\
\rowcolor{yellow!10}
\texttt{<information>} [St. Louis Cardinals, sports.sports\_team.arena\_stadium, Busch Stadium], [St. Louis Cardinals, sports.sports\_team.arena\_stadium, Roger Dean Stadium], [St. Louis Cardinals, sports.sports\_team.location, St. Louis], ... \texttt{</information>} \\
\rowcolor{blue!10}
\texttt{<think>} The query results indicate that the St. Louis Cardinals have multiple stadiums associated with them, including Busch Stadium and Sportsman's Park. Since Busch Stadium is mentioned multiple times and is the most recent name used, it is likely the primary home stadium of the team. Therefore, I will provide Busch Stadium as the answer. \texttt{</think>} \\
\hline
\rowcolor{green!20}
\texttt{<answer>} ["Busch Stadium"] \texttt{</answer>} \textbf{(EM=1.0 - correct answer with shortest path)} \\
\bottomrule
\end{longtable}
}



